% English translation of checked_pages/esq_119_042.tex; source: cote 119, site 42, PDF page 43, typescript page 40.
% AI-translated by GPT-6 Astra (OpenAI, ChatGPT Pro mode).
\scanpage{119/42}
\mspage{40}

of acquired habits, but tenacious, meticulous work,
probably over a long period, yielding eloquent harvests
day by day.

I would like to say a few more words about an earlier
reflection (late 60s ?), very close to the one just
discussed, inspired by Nash's ideas, which had
greatly struck me. Here, instead of defining axiomatically a notion
of a ``tame theory'' by giving a notion of a ``tame subset
of $\mathbb{R}^n$ '' satisfying certain conditions (chiefly of stability),
what I was after was an axiomatization of the notion of a ``smooth manifold'' and of the
differentiable formalism on such manifolds,
by giving, for each natural number $n$, a subring
$\unsure{\sourceglyph{119A}}_n$ of the ring of germs of real functions at the origin in $\mathbb{R}^n$.\tnote{TN-042-01}{The visible typed symbol used for this system of rings is reproduced from the scan. Its alphabet remains unresolved; the doubt marks preserve that uncertainty without assigning a letter family.}
These are the functions that will be admitted for expressing the ``changes
of chart'' for the corresponding notion of a $\unsure{\sourceglyph{119A}}$-manifold, and the
first task was to draw out\frapparatus{V-FR-042-001} a system of axioms on the
system $\unsure{\sourceglyph{119A}} = (\unsure{\sourceglyph{119A}}_n)_{n\in\underline{N}}$ that ensures this notion of a manifold a
flexibility comparable to that of a $C^\infty$, or real analytic
(or Nash), manifold. Depending on the kind of familiar constructions that
one wants to be able to carry out in the context of $\unsure{\sourceglyph{119A}}$-manifolds, the
relevant system of axioms is more or less restricted, or rich.
Very little suffices when it is only a matter of developing the
differential formalism, with the construction of jet bundles, De Rham
complexes, etc. To obtain a statement of the type ``quasi-finite implies
finite'' (for a map in a neighborhood of a point), which
emerged as a key statement in the local theory of analytic
spaces, one needs a stability axiom of a more delicate nature,
in Weierstrass's ``Vorbereitungssatz''\tnote{TN-042-02}{The German term ``Vorbereitungssatz'' means ``preparation theorem''.} (*). In other

\rule{4cm}{0.4pt}

(*) It may seem simpler to say that the (local) rings
$\unsure{\sourceglyph{119A}}_n$ are \uline{henselian}, which is equivalent. But in this latter
form it is by no means clear a priori that the condition in
question is in the nature of a stability condition, an
important circumstance, as will emerge in the reflections that follow.

\margin{42}
